\section{Related Work}
\label{sec:related}

\paragraph{Individual differential privacy.}
Classical $(\eps,\delta)$-DP~\cite{dwork2006calibrating} applies uniformly to
every record. Personalised DP, introduced by
\citet{jorgensen2015conservative}, allows per-record budgets and was first
realised by attribute suppression and weighting. The sampling-based
realisation of \citet{boenisch2023have} is the dominant deep-learning
implementation today and is the realisation under which the privacy
externality arises. \citet{feldman2021individual} pursue an alternative
\emph{filtering-based} iDP: records whose cumulative budget would be exceeded
in the next step are dropped from training. Filtering does not exhibit the
externality (each record's filter is independent), but in deep learning it
is impractical due to the catastrophic-forgetting risk reported by
\citet{boenisch2023have}. Our PEG is specific to the sampling-based
realisation; the externality vanishes in the sensitivity-based realisation,
which preserves the noise-to-sensitivity ratio across the dataset.

\paragraph{Excess privacy risk and mechanism comparison.}
Mechanisms calibrated to the same $(\eps,\delta)$ can have very different
privacy profiles and thus different realised
risks~\cite{kaissis2024beyond, balle2020privacy, dong2022gaussian}. The
$\Delta$-divergence of \citet{kaissis2024beyond} provides a notion of
approximate Blackwell-dominance that quantifies the worst-case excess MIA
advantage between mechanisms. \citet{kaiser2026your} use this framework to
motivate an $(\eps_i, \delta, \Delta)$-iDP contract that bounds excess
vulnerability and demonstrate empirically that without such a bound,
adversarial budget choices inflate the MIA susceptibility of $62\%$ of
targeted records. Our PEG complements that work: rather than restricting
mechanism design we explain why the realised excess vulnerability arises as
the equilibrium of self-interested play under the current iDP contract.

\paragraph{Game-theoretic privacy in collaborative learning.}
\citet{pejo2019together} introduced a two-player game on collaborative
learning, with players choosing privacy parameters $p \in [0, 1]$ in an
arbitrary privacy mechanism. They define the \emph{Price of Privacy} as the
relative accuracy loss at the equilibrium versus no privacy. Their
privacy mechanism's effect on accuracy depends on \emph{both} players'
choices, but its effect on \emph{privacy} depends only on the player's own
choice---there is no privacy externality. Our PEG transplants their
game-theoretic methodology into a setting where the externality is in the
privacy direction, not the accuracy direction. We retain the qualitative
contributions of their work (two-player setting, choice of $\eps$,
notion of equilibrium analysis) and extend them to the externality regime.

Earlier work studied related game-theoretic privacy questions in different
settings. \citet{ioannidis2013linear} model linear regression with private
features as a non-cooperative game. \citet{chessa2015game} study non-monetary
incentives in data analytics. \citet{pawlick2016stackelberg} model the
machine-learner-vs-data-perturber conflict as a Stackelberg game. None of
these treat the privacy externality of a shared sampling mechanism, which is
specific to the sampling-based realisation of iDP.

\paragraph{Membership inference attacks.}
The MIA advantage as we use it follows the formalisation of
\citet{carlini2022membership}, who introduced the likelihood-ratio attack
(LiRA) as a tight empirical estimator. \citet{carlini2022privacy} document
the privacy onion effect: not every data point is equally vulnerable at
the same nominal budget. Our framework predicts the \emph{average} realised
MIA advantage per group and abstracts over per-record onion heterogeneity.

\paragraph{Federated-DP regimes and the locus of the externality.}
It is worth being explicit about what the PEG is a game over.
\emph{The PEG is a game about budget coordination.} It applies to any
mechanism in which a coordinator solves for a shared noise multiplier
$\sigma$ and per-record (or per-group) sampling rates $p_g$ from a
public budget vector $\boldsymbol{\eps}$. The externality lemma
(\cref{lem:externality}) is a statement about that joint solve; the
equilibrium structure of \cref{sec:experiments} follows. Wherever a
coordinated mechanism instance of that shape appears, so does the
externality; wherever such a mechanism does not appear, the game we
analyse does not.

The sampling-based iDP of \citet{boenisch2023have} is one instantiation
of this abstract condition, and the one we build the paper on. Other
instantiations arise wherever heterogeneous per-record or per-silo
budgets need to be reconciled through a single central mechanism. A
common framing of DP in federated learning distinguishes three trust
regimes~\cite{kairouz2021advances}. Under \emph{central DP} the
aggregator is trusted and perturbs the aggregate on behalf of the
clients~\cite{geyer2017differentially}; under \emph{local DP} the
aggregator is untrusted and each client perturbs its own update before
transmission~\cite{truex2020ldp}; distributed DP bridges the two via
cryptographic intermediates such as secure
aggregation~\cite{bonawitz2017practical}. These trust regimes differ in
where the noise is added, not in how the budgets are coordinated. The
PEG applies to any of them that also carries a heterogeneous-budget
coordinator---most cleanly to central DP with per-silo or per-record
personalised budgets, exactly the regime under which the cross-silo
personalised-DP framework of \citet{liu2024crosssilo} operates. It
inherits, we conjecture, to distributed DP via secure aggregation,
because the coordinator still solves for $(\sigma, p_g)$ from a public
budget vector; the cryptography hides the gradients, not the budgets.
It does \emph{not} apply to pure local DP, because there is no shared
$\sigma$ to couple the players' sampling rates: each client's noise is
independent. Nor does it apply to alternative iDP realisations that
avoid the joint solve---sensitivity-based iDP~\cite{boenisch2023have}
preserves the noise-to-sensitivity ratio and so has no
externality, and filter-based iDP~\cite{feldman2021individual} enforces
per-record budgets by \emph{dropping} records that would exceed them
rather than coordinating them, again removing the coupling. In short,
the abstract sufficient condition is a coordinator jointly solving for
mechanism parameters from a public budget vector; any deployment that
matches that condition inherits the equilibrium structure of the PEG.

\paragraph{Federated learning incentive design.}
A separate literature studies privacy-aware incentive design in federated
learning: a server commits to incentives, and heterogeneous clients respond
with their own privacy choices, often modelled as a Stackelberg
game~\cite{li2024fedpcs, donahue2021optimality, tu2025idhfl, ma2024evolutionary}.
Our framework provides a foundational baseline on which such incentive
schemes must rest: it characterises the equilibrium that emerges when no
incentive is offered, which is what any incentive design needs to improve
upon.

\paragraph{Utility-noise tradeoff in DP-SGD.}
We motivate the functional form of the model-utility value $V(\sigma)$ via
the excess-empirical-risk bound of \citet{bassily2014private} for convex,
Lipschitz DP-ERM, who prove that the excess risk of DP-SGD scales as
$O(d / (n\eps)^2)$ at fixed $\delta$, equivalently as $O(\sigma^2 / n^2)$
under the moments accountant of \citet{abadi2016deep}. The convex setting
is restrictive; for deep learning, \citet{tramer2021private} provide
extensive empirical evidence that test accuracy degrades monotonically with
$\sigma$ on standard benchmarks. We use $V(\sigma) = 1/(1+\sigma^2)$ as a
directional proxy; the qualitative results of our paper do not depend on
the precise functional form.
